\scanpage{135}

% L1–6
$S$ affine, et si $(E_{10},\ldots,E_{\nu0})$ est un élément de $\Delta_H(E_0,q_0)$, \unsure{ceci|cela|celui-ci} se remonte \inserted{en un $(E_1,\ldots,E_\nu)\in\Delta_H(E,q)$.} Il suffit bien sûr (\deleted{\unsure{par|par la}}\inserted{utilisant} les propriétés de lissité de $\underline{\mathrm{Drap}}_H\,E\longrightarrow\underline{\mathrm{Grass}}_{i_\nu}(E)$) de prouver que $F_0=E_{\nu0}$ se remonte en un sous-Module $F=E_\nu$ \deleted{\unsure{loc.}} fact. direct et \uline{tot. isotrope}.

% L6–8, cancelled working
\deleted{On procède par récurrence sur $d=\operatorname{rg}F$. Si $d=1$ on commence à remonter \unsure{en $F_0$} \unsure{en} $F$}

% L9–13, cancelled working
\deleted{($\underline{\mathrm{Grass}}_{i_\nu}(E)$ lisse sur $S$) facteur direct, mais $q\mid F$\deleted{\unsure{$0$}} est à valeurs dans $\mathcal J$, \inserted{donc peut s'identifier à un él. de $\mathrm{Quad}(F_0,\mathcal J)$} pas néc. nul, \deleted{On a}\inserted{Or} une ``correction'' sur le choix de $F$ dépend de}

% L14, cancelled display
\[
 \deleted{\xi\in\operatorname{Hom}(F_0,E_0/F_0\otimes\mathcal J)\simeq
 \Gamma\bigl(\underline{\mathrm{Hom}}(F_0,E_0/F_0)\otimes\mathcal J\bigr)}
\]

% L15–18, cancelled working
\deleted{puis supposons que $F_0$ a une base $(e_{i0})_{1\leq i\leq d}$ ($d=i_\nu$). On \deleted{\unsure{montre}}\inserted{note} que chaque $e_{i0}$ peut se remonter en $e_i$ \inserted{$\in\Gamma E$ tel que $q(e_i)=0$,} (lissité de la quadrique projective définie par $q$), soit}

% L19–23, cancelled working
\deleted{\deleted{$F$ le module engendré par les} c'est la lissité de la quadrique projective. Si $d>1$, \unsure{OPS} $F_0=G_0+\mathcal O_S e_0$, $G_0$ de rang $d-1$, $e$ section de $F_0$. On remonte $G_0$ en $G\in\Delta_{\{d-1\}}(E,q)$, puis il reste à remonter $e_0$ en section $e$ de $E$}

% lower left
% archive sheet number: 134
